\chapter{Options and tuning}\label{ch:options}

\section{The Options dialog}\label{sec:options}

\menu{Tools $\rightarrow$ Options\ldots} opens the dialog of figure~\ref{fig:dialog-options}. The settings are saved when the
program closes, in the settings of the system (chapter~\ref{ch:install}).

\dialog{dialog-options}{The \menu{Options} dialog.}{9.5cm}

\subsection{General}

\begin{description}[style=nextline,leftmargin=1.2cm]
  \item[Interface size (in \%)] The size of the drawing of the window, from 100\% to 300\%. Only whole multiples of 100\% keep
        the pixels exact. The first start chooses it from the size of the screen (section~\ref{sec:window}).
  \item[Track line highlight step] The two numbers of the \texttt{HIGHLIGHT} field of the info area: every how many lines the lines of
        the tracks are highlighted (for example 8 and 4).
  \item[Use German notation] The note \texttt{B} is written \texttt{H}, as in the German notation.
  \item[Alternative track line numbering] The line numbers follow the highlight step.
  \item[Display accidentals as flats instead of sharps] \texttt{D\#} is written as \texttt{Eb}.
  \item[NTSC system speed (60 Hz)] The video system of the song: 60 frames a second instead of the 50 of PAL. The same switch is
        \key{CONTROL+F12}.
  \item[Audio buffer] How much sound the output of the system takes at a time, from 5 to 40 milliseconds. A small buffer
        makes the sound follow the keys and the MIDI keyboard sooner; a larger one keeps the sound clean when the computer or
        the sound server is busy. The default is 20\,ms on Linux, where the sound servers (PipeWire, PulseAudio) work in
        pieces of about that size, and 5\,ms on Windows and macOS. If the sound crackles or is distorted, choose 30 or 40\,ms.
  \item[Smooth scroll during playback] The lines of the tracks scroll one pixel at a time while a song plays.
  \item[Debug display] A line of internal values in the status bar. Enable it only if you know what it shows.
  \item[RMT Driver Version] The tracker driver, the player routine that makes the sound while you work (chapter~\ref{ch:drivers}). The
        driver must be the one the song was made with; the default is \emph{RMT 1.34 Patch16 by VinsCool}.
\end{description}

\subsection{Keyboard}

\begin{description}[style=nextline,leftmargin=1.2cm]
  \item[Keyboard layout] QWERTY, QWERTZ or AZERTY: where the piano of the note keys is. The first start picks the layout from the
        language of the keyboard.
  \item[Move to previous/next song line while track boundaries are crossed] The cursor passes from the end of a track to the next
        song line, and the other way round.
  \item[Remember octaves and volumes separately for each instrument] Each instrument keeps the octave and the volume you last used with it.
  \item[Reset the Atari sound routines each time ESC is pressed] \key{ESC} also starts the player routine again.
  \item[Prompt a save dialog box each time CTRL+S is pressed] \key{CONTROL+S} asks where to save, also from the menu and the
        toolbar.
\end{description}

\subsection{MIDI}

The MIDI input device, the touch response, the Atari volume offset and \emph{Record note off}: they are described in
chapter~\ref{ch:midi}.

\subsection{Paths}

The \menu{Paths\ldots} button opens a dialog with the default folders of the songs, the instruments and the tracks, each with a
\menu{Browse} button. The dialogs to open and save start there.

\section{The tuning}

\menu{Tools $\rightarrow$ Tuning\ldots} (figure~\ref{fig:dialog-tuning}) is a calculator for the pitch of the notes. The POKEY makes its
frequencies by dividing a clock, so not every pitch exists: the program computes the divisors that come nearest to the notes of
the temperament you choose, and the tracker plays with them.

\dialog{dialog-tuning}{The \menu{Tuning} dialog.}{10cm}

\begin{description}[style=nextline,leftmargin=1.2cm]
  \item[Base tuning (Hz) and Base note] The frequency of a reference note, by default \texttt{A-} at 440.8375\ldots\,Hz. It can be any
        number: \texttt{432}, \texttt{443.9}\ldots
  \item[Temperament] \emph{Equal Temperament} (the default) or another. The \emph{Ratio} table lists the interval of each step over the unison
        as a fraction.
  \item[Reset] puts the default values back; \menu{Test now} tries the tuning.
\end{description}

The tuning is saved with the configuration and is also the one a script uses when it exports a song (chapter~\ref{ch:scripting}).
